%% kapa: exam=1 round=22 year=2011 period=1 form=A subject=civil_law subject_ko=민법 seq=9 number=9
%% source: Q-Net 감정평가사 기출문제 (cst00309) artlSeq=5036790 file=extracted/1008/5036790/2011년 제22회 감정평가사 1차시험 기출문제/1교시.hwp
%% conversion: kordoc 4.24.51 -> tools/gen1.py (auto); review=proofread-sonnet; flags=none
\begin{question}{9}
조건과 기한에 관한 설명으로 옳지 않은 것은? (다툼이 있으면 판례에 의함)
\begin{choices}
  \item 기한의 도래는 소급효가 있다.
  \item 현상광고에서 정한 행위의 완료에는 조건이나 기한을 붙일 수 있다.
  \item 기성조건이 정지조건이면 조건 없는 법률행위가 되고, 해제조건이면 그 법률행위는 무효이다.
  \item 조건을 붙이는 것이 허용되지 않는 법률행위에 조건을 붙인 경우, 그 조건만을 분리하여 무효로 할 수 없다.
  \item 기한의 이익은 당사자의 특약이나 법률행위의 성질상 분명하지 않으면, 채무자에게 있는 것으로 추정된다.
\end{choices}
\end{question}
\answer{1}
